\setcounter{equation}{0}
\section{Conclusion}

Until quantum computers with enough qubits and low enough error rates to run useful applications become available, classical simulators will remain the primary venue for the development and analysis of quantum algorithms. The fundamental difficulty in simulating quantum systems on classical hardware --- the exponential growth of the state vector with the number of qubits --- motivates the use of parallel computation, and the simulation of small-to-medium registers across distributed clusters is a continuing area of active research.

This thesis began by introducing the four postulates of quantum mechanics, then built up the mathematical machinery (Dirac notation, the tensor product, unitary evolution, measurement) needed to describe quantum computation. It surveyed the historical landmarks of the field --- Benioff's quantum-mechanical Hamiltonian model, Feynman's quantum simulator proposal, Deutsch's quantum Turing machine, Shor's factoring algorithm, Grover's search algorithm --- and presented the three canonical algorithms in enough detail to be implementable: the Quantum Fourier Transform, Grover's unstructured search, and Shor's order-finding subroutine.

The implementation discussion was the part of the thesis that needed the most revision. The original strategy of distributing rows of a $2^n \times 2^n$ dense program matrix across MPI nodes was, as Section~\ref{sec:parallel} explains, the wrong primitive: the operator itself exceeds any conceivable memory budget for $n \gtrsim 20$, so no parallelisation strategy can rescue it. The right primitive is in-place sparse gate application directly on the state vector, with single-qubit gates costing $\Theta(2^n)$ work rather than $\Theta(2^{2n})$, and with distributed simulation arranged so that gates on low qubits cost no communication while gates on high qubits induce one pairwise exchange. The bundled 2004 C source remains as a historical record of the dense approach; a full implementation of the corrected library is left as future work.

A complete library along the lines sketched in Section~\ref{sec:parallel} and the Library Documentation chapter, with the QFT, Grover, and Shor circuits exposed as primitives, would prove useful for research into quantum algorithms --- and for teaching the subject without first requiring access to physical hardware. The Fourier transform and the algorithms that depend on it, in particular, are no longer optional features but the load-bearing core of the field.
